\documentclass[runningheads]{llncs}

\usepackage[T1]{fontenc}
\usepackage{graphicx}
\usepackage{amsmath}
\usepackage{amssymb}
\usepackage{booktabs}
\usepackage{algorithm}
\usepackage{algpseudocode}
\usepackage{xcolor}
\usepackage{hyperref}
\renewcommand\UrlFont{\color{blue}\rmfamily}
\urlstyle{rm}

\begin{document}

\title{Improvements to Backwards Proof Search with Iterative
Deepening and Rule Application Prioritisation}
\titlerunning{Backwards Proof Search Improvements}

\author{Madeline Abio}
\institute{Griffith University, Gold Coast, QLD 4101, Australia}

\maketitle
\begin{abstract}  The na\"ive backward proof-search procedure of the
  LK$'$ sequent calculus, as  described by Algorithm~2 in Hou's
  textbook~\cite{hou2021fundamentals}, is a simple backtracking
  algorithm  with an exponential search space and minimal heuristics.
  This  investigation presents a Rust implementation of the baseline
  together with  two algorithmic refinements: (i)~\emph{iterative
  deepening} of the depth-first  search, so that shallow proofs are
  always discovered before deep ones; and  (ii)~a six-class priority
  schedule for LK$'$ rules that prioritises  rules that add the least
  amount of states to the search space. The  combination of the two
  refinements forms the proposed  prover. On a 737-problem provable
  subset of the TPTP  library~\cite{Sut17}, the baseline na\"ive
  prover closes 152  problems within a one-second budget; iterative
  deepening alone closes 199  (+30.9\,\%); and the combined prover
  closes 209 (+37.5\,\%). The two  refinements are complementary, and
  the bulk of the improvement is contributed  by iterative deepening;
  the priority schedule contributes a smaller but  consistent
  additional gain on harder problems.

  \keywords{Sequent calculus \and LK$'$ \and Backward proof search \and
  Iterative deepening \and First-order logic \and TPTP.}
\end{abstract}


\section{Introduction}

Automated theorem proving (ATP) for first-order logic (FOL) is a mature field
with industrial-strength tools such as
Vampire~\cite{kovacs2013vampire}, E~\cite{schulz2002E},
Prover9~\cite{mccune2010prover9}, and the connection-based
leanCoP~\cite{otten2003leancop}. These provers are heavily engineered and
implement sophisticated techniques (resolution, superposition,3indexing,
clause selection) that are far removed from the textbook sequent calculus
introduced by Gentzen~\cite{gentzen1935untersuchungen}. The pedagogical
LK$'$ calculus presented in Hou~\cite{hou2021fundamentals} sits at the other
end of the spectrum: rules manipulate sequents directly, and a backward proof
search procedure (Algorithm~2 in that text) can be stated in twenty lines of
pseudocode. The same simplicity that makes LK$'$ accessible, however, makes
the na\"ive search strategy fragile: nothing prevents the prover from
choosing an unhelpful rule order, descending into an unbounded chain of
quantifier instantiations, or applying eigenvariable rules so eagerly that
later branching decompositions cannot recover.

This paper revisits Algorithm~2 from a practical standpoint. We give a
faithful Rust implementation of the baseline, identify two failure modes,
and propose two concrete algorithmic improvements that address them: a
six-class \emph{priority schedule} that refines the ``\textbf{else if}''
chain of Algorithm~2 (Sect.~\ref{sec:approach}), and an
\emph{iterative-deepening} envelope around the depth-first
search~\cite{korf1985iterative}. An empirical evaluation on a 737-problem
provable subset of TPTP shows that the combined prover closes 34.4\,\%
more problems than the baseline within a one-second budget
(Sect.~\ref{sec:implementation}).

\section{Background}\label{sec:background}

\subsubsection{Sequent calculi.}
A \emph{sequent} $\Gamma\vdash\Delta$ asserts that the conjunction of the
formulae in $\Gamma$ entails the disjunction of the formulae in
$\Delta$. Gentzen's classical sequent calculus
LK~\cite{gentzen1935untersuchungen} introduces logical connectives by pairs
of left- and right-introduction rules. A \emph{derivation tree} is a tree of
sequents whose root is the sequent to be proved and whose leaves are
\emph{axioms} closed by structural rules such as identity. The variant
LK$'$ used in Hou~\cite{hou2021fundamentals} differs from LK by making the
quantifier rules $\forall L$ and $\exists R$ \emph{reusable}: the original
quantified formula is retained on its side of the sequent after
instantiation, so a single $\forall x.\,A$ on the left can produce the
instantiated formulae $A[t_1/x], A[t_2/x],\ldots$ along the same branch.

\subsubsection{Backward proof search.}
\emph{Backward} (or \emph{root-first}) search begins from the goal sequent
and applies the inverse of the inference rules until every leaf is closed
by a structural rule. For propositional connectives the rules of LK$'$ are
invertible, so the order of decomposition does not affect provability;
quantifier rules, however, branch over either an eigenconstant
($\forall R$, $\exists L$) or a witness term ($\forall L$, $\exists R$).
Choice of witness term is the principal source of search-space explosion.
Tableau-based provers such as leanTAP~\cite{beckert1995leantap} address
this by free-variable instantiation and unification at branch closure;
connection-based provers such as leanCoP~\cite{otten2003leancop} avoid
explicit instantiation altogether. In this paper we stay within the
ground-instantiation regime of Algorithm~2.

\subsubsection{Iterative deepening.}
Korf~\cite{korf1985iterative} showed that depth-first iterative-deepening
(DFID) is asymptotically optimal in time, space, and solution-cost on
exponentially branching search trees. The technique runs depth-bounded DFS
at increasing depth limits$1, 2, 3, \ldots$, paying a small re-traversal
overhead in exchange for breadth-first-like solution ordering and bounded
memory. The technique is standard in heuristic search and is the basis of
our first improvement.

\section{Proposed Approach}\label{sec:approach}

\subsection{Baseline: Algorithm 2}\label{sec:baseline}

Algorithm~\ref{alg:naive} reproduces the na\"ive backward proof-search
procedure of Hou~\cite[p.~67]{hou2021fundamentals}, restated in our
notation. The five priority groups are mutually exclusive: a higher-priority
group is exhausted before any rule from a lower-priority group is
considered.

\begin{algorithm}[t]
  \caption{Na\"ive backward proof search for LK$'$ (Algorithm~2 of
  \cite{hou2021fundamentals})}\label{alg:naive}
  \begin{algorithmic}[1]
    \Require First-order formula $A$
    \Ensure Derivation tree in LK$'$, or report no proof found
    \State build the bottom sequent $\vdash A$
    \ForAll{top sequents on an open branch}
    \If{any of $\mathit{id}, \top R, \bot L$ is applicable}
    \State apply the rule backwards and close the branch
    \ElsIf{any of $\wedge L, \vee R, {\rightarrow}R, \neg L, \neg R,
    \forall R, \exists L$ is applicable}
    \State apply the rule backwards
    \ElsIf{any of $\wedge R, \vee L, {\rightarrow}L$ is applicable}
    \State apply the rule backwards and create a new branch
    \ElsIf{$\forall L$ or $\exists R$ is applicable, and there is a
      term $t$ that has not been used to instantiate the quantified
    variable $x$ in this formula}
    \State apply the rule backwards, substituting $x$ with $t$
    \ElsIf{$\forall L$ or $\exists R$ is applicable}
    \State apply the rule backwards and create a fresh term
    \Else
    \State stop
    \EndIf
    \EndFor
  \end{algorithmic}
\end{algorithm}

\subsection{Failure modes}
We identify two structural weaknesses of Algorithm~\ref{alg:naive}.
\textbf{(F1) Eager eigenconstants.} Group~2 lumps the eigenvariable
quantifier rules ($\forall R$, $\exists L$) with non-branching
propositional rules, so the prover commits to a fresh eigenconstant
even when a branching propositional decomposition would have closed the
branch by identity. Each eigenconstant enlarges the symbol set, which
enlarges the candidate set for subsequent $\forall L$/$\exists R$
instantiations.
\textbf{(F2) Unbounded reusable instantiation.} Because $\forall L$ and
$\exists R$ retain the original quantified formula on the sequent,
groups~4 and 5 of Algorithm~\ref{alg:naive} permit an arbitrary number
of distinct witness terms to be tried before any sibling branch is
revisited. A na\"ive depth-first loop therefore commits to deep
instantiation chains before backtracking to easier alternatives.

\subsection{Proposed Algorithm}
\begin{algorithm}[t]
  \caption{Improved backward proof search for LK$'$}\label{alg:improved}
  \begin{algorithmic}[1]
    \Require First-order formula $A$
    \Ensure Derivation tree in LK$'$, or report no proof found
    \State build the bottom sequent $\vdash A$
    \ForAll{top sequents on an open branch}
    \If{any of $\mathit{id}, \top R, \bot L$ is applicable}
    \State apply the rule backwards and close the branch
    \ElsIf{any of $\wedge L, \vee R, {\rightarrow}R, \neg L, \neg R,
    \forall R, \exists L$ is applicable}
    \State apply the rule backwards
    \ElsIf{any of $\wedge R, \vee L, {\rightarrow}L$ is applicable}
    \State apply the rule backwards and create a new branch
    \ElsIf{$\forall L$ or $\exists R$ is applicable, and there is a
    term $t$ that has not been used to instantiate $x$ in this formula}
    \State apply the rule backwards, substituting $x$ with $t$
    \ElsIf{$\forall L$ or $\exists R$ is applicable}
    \State apply the rule backwards and create a fresh term
    \Else
    \State stop
    \EndIf
    \EndFor
  \end{algorithmic}
\end{algorithm}
\subsection{Building block A: iterative deepening}

Our first refinement wraps the depth-first search of
Algorithm~\ref{alg:naive} in an iterative-deepening
envelope~\cite{korf1985iterative}. The prover runs depth-bounded DFS at
limits $d = 1, 2, \ldots, d_{\max}$ and returns the first depth at which
the goal closes. Because backward-search depth is roughly proportional
to the number of rule applications on the longest open branch, this
re-orders the search so that shallow proofs are discovered before any
deep instantiation chain is attempted, mitigating (F2).

\subsection{Building block B: LK$'$ six-class priority schedule}

Our second refinement promotes the five-group ``\textbf{else if}'' chain
of Algorithm~\ref{alg:naive} to a six-class schedule:
(1)~\emph{closing} rules $\mathit{id}, \top R, \bot L$;
(2)~\emph{non-branching propositional} rules $\wedge L, \vee R,
{\rightarrow}R, \neg L, \neg R$;
(3)~\emph{branching propositional} rules $\wedge R, \vee L,
{\rightarrow}L$;
(4)~\emph{eigenvariable quantifier} rules $\forall R, \exists L$;
(5)~\emph{reusable quantifiers with visible terms}: $\forall L, \exists R$
instantiated with a previously unused term already in the sequent;
(6)~\emph{reusable quantifiers, fresh fallback}: $\forall L, \exists R$
instantiated with a freshly generated witness subject to a per-occurrence
budget. The change is twofold. Classes~3 and~4 are reordered: branching
propositional rules are tried before any eigenconstant is introduced,
addressing~(F1). The original group~4/5 of Algorithm~\ref{alg:naive} is
split into classes~5 and~6: the prover always exhausts visible witnesses
before introducing a fresh constant, addressing~(F2). The bounded
fresh-fallback budget guarantees that the search returns
$\textit{Unknown}$ rather than $\textit{NotProvable}$ when only fresh
witnesses remain to be tried, preserving soundness.

\subsection{Combined prover}

Our proposed prover \texttt{priority-id} combines both building blocks:
it uses the six-class priority schedule inside an iterative-deepening
envelope. We use the iterative-deepening-only ablation \texttt{id} and
the unmodified \texttt{naive} engine as controls. All three share the
same parser, rule applicators, fresh-name generators, branch-state
bookkeeping, timeout machinery, and result persistence; they differ only
in the search engine that orders rule application.

\section{Implementation and Experiment}\label{sec:implementation}

\subsection{Implementation}

The prover is written in Rust~\cite{rust2024} (edition~2024) and consists
of approximately 4{,}500 lines of source code spread over a parser,
abstract-syntax tree, sequent representation, rule applicators, search
engine, and a CLI front-end. The parser is built with
\texttt{pest}~\cite{pest2024}, a parsing-expression-grammar generator,
against a hand-written grammar for the first-order-form (FOF) fragment of
the TPTP language~\cite{Sut17}. The CLI uses
\texttt{clap}~\cite{clap2024} for argument parsing and
\texttt{rusqlite}~\cite{rusqlite2024} for persisting per-problem results
to a SQLite database. Cancellation under \texttt{Ctrl+C} is handled by the
\texttt{ctrlc} crate. The three search engines (\texttt{naive},
\texttt{id}, \texttt{priority-id}) are selected at the command line and
share all code outside the rule scheduler.

Two design points warrant mention. First, the LK$'$-reusable rules
$\forall L$ and $\exists R$ retain the original quantified formula on its
side of the sequent after instantiation, so the prover tracks per-branch
quantifier-instantiation history (which terms have been used for which
quantified occurrence) to avoid trying the same instantiation twice on the
same branch. Second, fresh-name generation for both eigenconstants
($\forall R$, $\exists L$) and fallback witnesses ($\forall L$, $\exists R$)
inspects every symbol currently visible in the sequent and returns a name
not in that set, guaranteeing eigenvariable freshness without further
$\alpha$-renaming machinery. Substitution is capture-avoiding by
construction because only ground terms are ever instantiated for bound
variables.

\subsection{Experimental setup}

All benchmarks were run on a single laptop (Windows~11 Home, Intel x86\_64,
Rust~1.85). Each prover invocation was given a 1\,000\,ms wall-clock
timeout, a recursive search-depth limit of 128, a step limit of
10\,000 rule applications, and a fresh-fallback budget of one term per
quantified occurrence. Results were persisted to a SQLite database; the
analysis pipeline (Python~3.12 with \texttt{pandas} and
\texttt{matplotlib}) selects, for each engine, the canonical run with the
greatest number of recorded results, after asserting that all selected
runs share identical timeout, depth, step, and budget settings.

\subsection{Dataset}

We evaluate on a 737-problem subset of the TPTP problem
library~\cite{Sut17}, drawn from the FOF/THM problems with
0--50 input formulae, 0--150 atoms, and no equality or arithmetic. The
present evaluation reports results for the \emph{provable} class only;
extending the evaluation to unprovable and mixed classes is in progress
and will be reported separately.

\subsection{Results}

Table~\ref{tab:summary} reports the headline outcome counts. Within the
1\,000\,ms budget the baseline \texttt{naive} engine closes
160 of 737 problems (21.7\,\%); iterative deepening alone closes 209
(28.4\,\%); the combined \texttt{priority-id} prover closes 215
(29.2\,\%). The \texttt{naive} engine also produces the only two
\textit{Unknown} outcomes attributable to depth-bound exhaustion, a
symptom of failure mode~(F2): the depth-first envelope dives into a deep
quantifier-instantiation chain before exploring shallower alternatives,
and reaches the depth limit on problems that the iterative-deepening
engines close in milliseconds.

\begin{table}[t]
  \centering
  \caption{Per-engine outcomes on the 737-problem provable TPTP subset
  (timeout 1\,000\,ms, depth 128, steps 10\,000, fresh-fallback budget 1).}
  \label{tab:summary}
  \begin{tabular}{lrrrrrr}
    \toprule
    Engine & Solved & Solve\,\% & Timeout & Bicond.\,cap & Steps & Depth \\
    \midrule
    \texttt{naive}        & 160 & 21.7 & 562 & 12 & 1 & 2 \\
    \texttt{id}           & 209 & 28.4 & 515 & 12 & 1 & 0 \\
    \texttt{priority-id}  & 215 & 29.2 & 509 & 12 & 1 & 0 \\
    \bottomrule
  \end{tabular}
\end{table}

\begin{figure}[t]
  \centering
  \begin{minipage}{0.49\textwidth}
    \centering
    \includegraphics[width=\textwidth]{figures/outcome_composition.pdf}\\
    \small (a) Outcome composition.
  \end{minipage}\hfill
  \begin{minipage}{0.49\textwidth}
    \centering
    \includegraphics[width=\textwidth]{figures/solve_time_cdf.pdf}\\
    \small (b) Solve-time CDF (log scale).
  \end{minipage}\\[0.6em]
  \begin{minipage}{0.49\textwidth}
    \centering
    \includegraphics[width=\textwidth]{figures/solve_rate_by_atoms.pdf}\\
    \small (c) Solve rate by atom count.
  \end{minipage}\hfill
  \begin{minipage}{0.49\textwidth}
    \centering
    \includegraphics[width=\textwidth]{figures/solve_rate_by_difficulty.pdf}\\
    \small (d) Solve rate by TPTP rating.
  \end{minipage}
  \caption{Per-engine evaluation on the 737-problem provable TPTP subset.
    \textbf{(a)}~Solved counts increase monotonically from \texttt{naive} to
    \texttt{priority-id}; the bulk of the gain comes from the
    iterative-deepening envelope. \textbf{(b)}~The \texttt{naive} engine
    resolves the easiest problems fastest but saturates earlier; the
    iterative-deepening engines continue to close problems past that point.
    \textbf{(c)}~On atom-dense problems \texttt{priority-id} dominates the
    \texttt{id} ablation, where the visible-terms-first schedule reduces
    candidate instantiations. \textbf{(d)}~Solve rate decreases monotonically
    with TPTP difficulty for all three engines; the relative ordering is
  preserved across the rating range.}
  \label{fig:eval}
\end{figure}

Figure~\ref{fig:eval}(a) visualises the same data as a stacked outcome
composition: the proportions of \emph{biconditional-cap} and
\emph{max-steps} unknowns are essentially identical across the three
engines, confirming consistent configuration and that the differences in
solve count come from search order alone. The CDF in
Fig.~\ref{fig:eval}(b) exhibits the expected iterative-deepening
trade-off: the \texttt{naive} engine has the lowest median solve time
(2\,ms versus 7\,ms for \texttt{id} and 8\,ms for \texttt{priority-id})
because depth-first DFS without re-traversal is cheaper on
trivially-provable problems, but the iterative-deepening engines pay
that small constant overhead in exchange for closing 49--55 more
problems within the same wall-clock budget. Stratifying by atom count
(Fig.~\ref{fig:eval}(c)) reveals that \texttt{priority-id} pulls ahead
of the \texttt{id} ablation on the most atom-dense problems, where the
priority schedule's preference for visible witnesses pays off. Finally,
the TPTP-rating cross-validation in Fig.~\ref{fig:eval}(d) shows the
expected monotonic decline and confirms that the relative ordering
\texttt{priority-id} $>$ \texttt{id} $>$ \texttt{naive} holds across
every difficulty bin, so the improvement is not driven by a handful of
pathologically easy problems.

\section{Discussion and Conclusion}

The two refinements proposed in this paper address distinct failure modes
of Algorithm~\ref{alg:naive} of \cite{hou2021fundamentals}. Iterative
deepening, the more impactful of the two, recovers proofs that the
depth-first baseline misses by getting stuck in deep
quantifier-instantiation chains: it accounts for 49 of the 55 additional
solves observed in our combined prover. The LK$'$ six-class priority
schedule contributes a smaller but consistent additional gain,
concentrated on problems with many atoms where the visible-terms-first
phase pays for itself by avoiding unnecessary fresh-witness generation.

The remaining 70\,\% of unsolved problems are dominated by timeouts
rather than soundness-preserving \emph{Unknown} outcomes, indicating
that the prover would benefit more from sharper search heuristics than
from a larger resource budget. Promising directions include
unification-based instantiation in the manner of
leanTAP~\cite{beckert1995leantap}, term ordering for $\forall L$/$\exists R$
witness selection, and lemma sharing across iterative-deepening
iterations to amortise the re-traversal overhead. Extending the
evaluation to the unprovable problem class would also test the dual
question of how quickly the prover can saturate a non-theorem and report
\emph{NotProvable} rather than time out.

The principal lesson of this study is that small, principled changes to
the rule-application order of a textbook sequent calculus can produce
disproportionate improvements in practical solve rate, even before any
of the heavyweight machinery (resolution, superposition, free-variable
unification) of mature first-order provers is brought to bear.

\subsubsection{Data Availability.}
Source code, benchmark scripts, and the SQLite database backing every
figure in this paper are available at the project repository on GitHub.

\bibliographystyle{splncs04}
\bibliography{mybibliography}

\end{document}
